\section{Analysis Layer~3 --- Biophilic Interaction Forms}
\label{sec:ch5-layer3}

This section presents six \emph{biophilic interaction forms} that articulate how occupants might interact with future biophilic indoor environments (\textbf{RQ3}). Expressed as verbs, they foreground interaction as experienced and enacted by people. Each form is described through (1) its interaction logic, (2) supporting system mechanisms, and (3) an example use case illustrating how it could occur in practice. An overview is provided in \autoref{tab:biophilic-interaction-forms}. The interaction forms are grounded in patterns identified across the empirical material. The accompanying system mechanisms are author-developed interpretive extrapolations included to illustrate how these interaction logics might be realised. They are therefore not direct outputs of the workshops, nor validated technical solutions, but speculative constructs intended to make the interaction logic tangible and discussable. Several mechanisms (e.g. variable-resistance flooring or distributed environmental modulation) would require substantial technical development and real-world evaluation.

\paragraph{\textbf{Anchoring}}
Anchoring foregrounds bodily contact with the ground as an active channel through which stability and presence are registered. Like walking on sand, occupants perceive an amplified bodily weight and downward pull through interior surfaces that respond to pressure and posture, making gravity perceptible rather than neutralised. 

\emph{Supporting Mechanism.} Pressure-sensitive flooring coupled with variable-resistance surface actuation detects stance duration, weight distribution, and gait change, and responds by locally modulating surface compliance and friction.

\emph{Use Case.} When an occupant enters a corridor after extended desk work and naturally slows their pace while transitioning between activities, floor sensors detect prolonged stance and reduced walking speeds. In response, the flooring increases resistance and dampens impact by mechanically yielding under load, slowing movement and intensifying contact with the ground. The occupant pauses briefly, using proprioceptive feedback rather than visual cues to register bodily presence before continuing onward.

\paragraph{\textbf{Synchronising}}
Synchronising is where bodily timing becomes coordinated with shared environmental rhythms occurring over time rather than through discrete commands. Similar to surfing small waves, synchronising describes interaction in which occupants align movement, pace, or stillness with recurring environmental rhythms, rather than acting independently of temporal patterns. 

\emph{Supporting Mechanism.} Pressure-sensitive flooring, ceiling-mounted presence sensors, and ambient environmental sensing such as CO$_2$ monitoring detect aggregate movement patterns and postural changes.

\emph{Use Case.} In a long corridor connecting workspaces to meeting rooms, presence sensors detect walking as occupants transition between focused work and collaborative activities. In response, overhead lighting and ambient sound begin cycling in a slow, regular sequence. As light brightens and dims and sound rises and falls, the occupant is implicitly nudged to slow down or pause briefly, allowing bodily rhythm to align with the environmental cycle. This supports smoother transitions between activities by allowing pace and attention to settle.

\paragraph{\textbf{Dwelling}}
Dwelling is conceptualised as habitation, where remaining in place becomes an active mode of interaction through which comfort and enclosure accumulate with duration. Like resting under a tree's shadow, dwelling describes interaction in which sustained presence leads to gradually reshaped bodily experience. 

\emph{Supporting Mechanism.} Pressure-sensitive seating and floor-based presence sensing detect prolonged presence and low movement variability over time.

\emph{Use Case.} In a semi-enclosed alcove adjacent to an open workspace, an occupant sits down after completing a demanding task. Sensors detect continuous seated presence with minimal movement. In response, sound masking gradually increases, lighting becomes more diffuse. As these conditions accumulate, the space increasingly supports cognitive recovery and reduced environmental awareness, allowing the occupant to remain without adjusting settings---rewarding duration.

\paragraph{\textbf{Inferring}}
Inferring is when occupants actively probe conditions and interpret indirect multisensory feedback rather than receiving explicit visual or textual information. This interaction enhances spatial awareness through embodied learning, extends accessibility by making non-visual wayfinding strategies available to all occupants, and supports visual rest from screens. Like wandering in a forest at dawn, inferring describes interaction in which occupants seek spatial or material conditions without relying primarily on vision, using non-visual cues such as sound, vibration, airflow, or thermal change. 

\emph{Supporting Mechanism.} Proximity sensing combined with vibration, airflow modulation, and directional sound provides graded multisensory feedback in contexts where visual information is intentionally reduced.

\emph{Use Case.} In an open-plan office with multiple adjacent work zones, visual boundaries are  minimal. Environmental sensing modulates non-visual cues instead: quieter zones produce softer acoustic textures and reduced floor vibration, while more active zones subtly increase low-frequency sound and air movement from ventilation diffusers. As occupants move through the space, they sense these gradients through their bodies rather than through signs or screens. By slowing, changing direction, or pausing, occupants infer where focused or social activity is taking place and choose where to settle. Over time, the office becomes legible through embodied interpretation rather than spatial coding.

\paragraph{\textbf{Withholding}}
Withholding describes a form of interaction in which the environment remains perceptible and affectively present, yet deliberately resists approach, touch, or prolonged engagement. Rather than inviting participation or occupation, these spaces support pass-by encounter and attentive distance, foregrounding presence without access. Withholding draws from natural situations that are experientially rich but not enterable---such as observing lightning from afar or sensing the pull of a cliff edge. Withholding is a form of anti-interaction: a way of being in relation to an environment without fully entering into it.

\emph{Supporting Mechanism.} Material and environmental cues such as sharp light, directional sound, or visually arresting but unreachable features shape proximity and attention. Smart sensing can be used to control these qualities, for example by reducing responsiveness as proximity increases rather than amplifying it.

\emph{Use Case.} A narrow transition corridor incorporates textured surfaces, directional lighting, and ambient natural sound that shift as people pass. The space offers no seating or touchable elements and does not reward lingering. While transitory, it leaves a strong sensory impression that punctuates everyday movement.

\begingroup
\fontsize{6.8}{8.1}\selectfont
\setlength{\tabcolsep}{2.1pt}
\renewcommand{\arraystretch}{1.04}
\setlength{\LTleft}{0pt}
\setlength{\LTright}{0pt}

\begin{longtable}{@{}%
  >{\RaggedRight\arraybackslash}p{.16\linewidth}
  >{\RaggedRight\arraybackslash}p{.25\linewidth}
  >{\RaggedRight\arraybackslash}p{.15\linewidth}
  >{\RaggedRight\arraybackslash}p{.23\linewidth}
  >{\RaggedRight\arraybackslash}p{.16\linewidth}@{}}

\caption{Overview of the biophilic interaction forms and their underlying interaction logic, analytically derived from the session data. Examples of bodily actions and spatial contexts are grounded in abstracted interaction episodes identified during analysis (cf.~\autoref{subsec:ch5-layer3}); the enabling system mechanisms are author-developed extrapolations included to make the interaction logic tangible.}
\label{tab:biophilic-interaction-forms}\\

\toprule
\textbf{Interaction Form} &
\textbf{Interaction Logic} &
\textbf{Bodily Actions Involved} &
\textbf{Example System Mechanisms} &
\textbf{Example Spatial Contexts} \\
\midrule
\endfirsthead

\caption[]{Overview of the biophilic interaction forms and their underlying interaction logic, continued.}\\
\toprule
\textbf{Interaction Form} &
\textbf{Interaction Logic} &
\textbf{Bodily Actions Involved} &
\textbf{Example System Mechanisms} &
\textbf{Example Spatial Contexts} \\
\midrule
\endhead

\midrule
\multicolumn{5}{r}{\footnotesize Continued on next page}\\
\endfoot

\bottomrule
\endlastfoot

Anchoring &
Bodily weight and gravity are made perceptible through responsive ground contact rather than neutralised by flat, passive floors. &
Slowing, weight-bearing, pausing &
Pressure-sensitive flooring with variable resistance or compliance responding to stance duration and gait change &
Corridors, transition zones, entry passages \\

\midrule

Synchronising &
Movement and stillness align with shared environmental rhythms rather than unfolding independently of temporal patterns. &
Walking, pacing, pausing &
Presence sensing, pressure-sensitive flooring, ambient sensing driving rhythmic modulation of light and sound &
Corridors between work areas, transitional passages \\

\midrule

Dwelling &
Sustained presence allows environmental conditions to accumulate over time rather than responding immediately on entry. &
Sitting, remaining, stillness &
Pressure-sensitive seating and floor-based presence sensing enabling gradual sound, light, and thermal adaptation &
Seating alcoves, semi-enclosed work or recovery spaces \\

\midrule

Inferring &
Spatial and material conditions are understood through non-visual cues rather than explicit visual legibility. &
Probing, testing, cautious stepping &
Proximity sensing, vibration, airflow modulation, and directional sound providing graded multisensory feedback &
Low-lit wellness rooms, sensory exploration areas, visually reduced spaces \\

\midrule

Withholding &
Environmental qualities remain perceptible but resist approach, lingering, or direct engagement, shaping interaction through maintained distance rather than invitation. &
Passing by, orienting without approaching, maintaining distance &
Environmental sensing combined with lighting, acoustic, or material actuation that reduces responsiveness as proximity increases &
Transitional corridors, threshold spaces, circulation zones adjacent to focal areas \\

\midrule

Accepting &
Environmental mediation withdraws temporarily, allowing occupants to experience unbuffered conditions through controlled absence. &
Remaining exposed, lingering, withholding adjustment &
Presence detection triggering gradual disengagement of sound masking, lighting control, and climate regulation &
Transitional alcoves, threshold spaces between regulated and less-regulated zones \\

\end{longtable}

\Description{A table summarising the biophilic interaction forms, including for each form the core interaction logic, examples of possible bodily actions, illustrative system mechanisms, and relevant spatial contexts.}

\endgroup




\paragraph{\textbf{Accepting}}
Accepting is conceptualised as exposure through absence of regulation, where occupants experience conditions without technological smoothing. Interaction with nature often involves accepting a level of discomfort. Technological mediation to optimise environmental comfort can progressively withdraw, allowing unbuffered light, sound, and temperature to become perceptible rather than immediately corrected. 

\emph{Supporting Mechanism.} Distributed indoor and outdoor environmental sensing is combined with kinetic actuation. Control logic prioritises continuity and variation over equilibrium.

\emph{Use Case.} In a waiting space near an atrium, overhead textile panels are actuated by a slow robotic system informed by outdoor wind and indoor air-pressure sensing. As conditions shift, panels open and fold, intermittently channelling cooler air and altering shadow patterns. The system does not respond to presence or dwell time, nor does it aim to maintain steady conditions. Occupants experience moments of draught, shifting light, and changing enclosure. Remaining in the space involves accepting these variations, while those seeking stability can move elsewhere. Technology remains visibly active, sustaining exposure rather than eliminating it.


